\section{Model DIA Legal Language}
\label{sec:legal}

\subsection{Purpose}

This section provides model statutory language for state legislatures
considering the \dia{}.
The language is designed to be directly adoptable and covers three elements:
(1) mandatory manifest filing with the Secretary of State,
(2) retention period, and (3) public verification rights.

\subsection{Model Statutory Language}

\noindent\textbf{Section [X]. Algorithmic Redistricting Manifest.}

\medskip
(a) \textit{Manifest generation.} Any redistricting plan adopted pursuant
to this Act shall be accompanied by a cryptographic computation manifest
(``Manifest'') generated by the certified redistricting algorithm.
The Manifest shall contain, at minimum:
\begin{enumerate}
\item The SHA-256 hash of the redistricting algorithm binary used to
  generate the plan, identifying the specific version of the algorithm;
\item The SHA-256 hash of the parameter configuration file used to
  generate the plan;
\item The SHA-256 hash of each Census Bureau data file used as input
  to the algorithm, as provided by the U.S.\ Census Bureau's official
  redistricting data release;
\item The SHA-256 hash of the adjacency graph constructed from Census
  Bureau TIGER/Line shapefiles;
\item The SHA-256 hash of the final district assignment output;
\item A self-referential integrity hash computed as the SHA-256 hash
  of the foregoing contents of the Manifest.
\end{enumerate}

\medskip
(b) \textit{Filing requirement.} The redistricting authority shall file
the Manifest with the Secretary of State within thirty (30) days of plan
adoption.
The Secretary of State shall make the Manifest available for public
inspection on the Secretary of State's official website within five (5)
business days of receipt.

\medskip
(c) \textit{Retention period.} The redistricting authority shall retain
all source data, parameter configuration files, algorithm binary, and
intermediate computation outputs referenced in the Manifest for a period
of not less than ten (10) years following the effective date of the
redistricting plan, or until such time as any pending litigation arising
from the plan is finally resolved, whichever is later.

\medskip
(d) \textit{Verification rights.} Any registered voter in the state or
any candidate who has filed for office in the state shall have the right
to:
\begin{enumerate}
\item Obtain a copy of the Manifest from the Secretary of State;
\item Independently verify the Manifest by running the certified
  redistricting algorithm's verification tool against the Manifest;
\item Request from the redistricting authority, at no charge, a copy
  of any source data or parameter file referenced in the Manifest.
\end{enumerate}

\medskip
(e) \textit{Violation.} A redistricting plan whose Manifest fails
verification under the certified algorithm's verification tool shall
be voidable at the instance of any person who has standing to challenge
the plan.
The redistricting authority bears the burden of demonstrating that a
verification failure is attributable to technical error rather than
data manipulation.

\medskip
(f) \textit{Expert testimony.} In any proceeding challenging a redistricting
plan adopted under this Act, the Manifest and the results of the
verification tool constitute admissible evidence of the computation's
provenance and integrity, subject to the provisions of the state rules
of evidence governing scientific and technical evidence.

\subsection{Commentary}

The 30-day filing deadline gives the redistricting authority time to
complete verification but prevents indefinite delay.
The 10-year retention period covers the full decade until the next
redistricting cycle plus a buffer for litigation that may extend beyond
the cycle.
The verification right for any registered voter reflects the public
nature of redistricting and ensures that opponents do not need to file
formal legal proceedings to access verification tools.

The burden-shifting in subsection (e) reflects the technical reality that
SHA-256 verification failures are almost never technical errors; the
probability of a false failure under correct computation is negligible.
Requiring the redistricting authority to explain a failure creates an
accountability mechanism without creating a presumption of fraud for
genuinely technical edge cases.
